\begin{lemma}[A zero-divider minimum hierarchy counterexample cannot have one tail addition]
Fix a finite loopless multigraph with its individual edge identities, a
palette of $k$ colors, and a distinguished edge $e_0$.  Let
$(T,\varphi_n)$ be a positive-rung extended Tashkinov tree constructed
from an arbitrary certified input coloring, let $\mathcal B$ be a
hierarchy base under $\varphi_n$, and let $C$ be a hierarchy
counterexample over that base.  Write $H=C.\mathit{hierarchy}$, let $q$
be its divider count, and put
\[
 A=\noderef{HierarchyStarredStage}(H,q),\qquad
 \ell=|T|-|A|,
\]
and let
\[
 T(y_{\ell-1})=
 \noderef{TreeSequencePrefix}
   \bigl(T,\min\{|A|+(\ell-1),|T|\}\bigr).
\]
Assume that $\varphi_n$ is proper and leaves exactly $e_0$ uncolored.

Assume, uniformly for every positive-rung maximum ETT on the same graph,
palette, and distinguished edge, with arbitrary certified input and
generating colorings and with rung number smaller than that of $T$, all
six lower-rung assertions of the technical theorem: elementarity;
terminal interchangeability; at every PE slot, unconditional stability
before the support, stability on the preceding ladder tree when the
first-appearance index is greater than one, and disjointness of stored
pairs at distinct PE slots with distinct supports and equal
first-appearance index; the terminal PE component assertion; the
defective-endpoint order assertion; and the mod-coloring assertion.
Assume that every same-rung maximum ETT with a good hierarchy having
fewer than $q$ dividers is elementary.

Assume that $C$ is lexicographically minimal among hierarchy
counterexamples whose ETTs have the same rung number as $T$ and whose
hierarchies have exactly $q$ dividers, allowing the certified input
coloring, generating coloring, ETT, base, and hierarchy to vary.  Thus
no such $C'$ has
\[
 \noderef{HierarchyCounterexampleLess}
   (\noderef{HierarchyCounterexampleMeasure}(C'),
    \noderef{HierarchyCounterexampleMeasure}(C)).
\]
Assume also the large-branch estimate
\[
 2n+11\leq |\overline\varphi_n(T_n)|,
\]
where $n$ is the rung number and $T_n$ is the terminal ladder tree.

If
\[
 \ell=|T|-|A|=1,\qquad
 V(T(y_{\ell-1}))\text{ is elementary under }\varphi_n,
 \qquad q=0,
\]
then a contradiction follows.
\end{lemma}
\begin{proof}
All ordered prefixes in this proof retain the individual edge elements
of the multigraph.  The zero-index clause of
\noderef{HierarchyStarredStage}, the `coreInsideFirst' and
terminal-stage fields of \noderef{TashkinovHierarchy}, and
$\ell=1$ give $A\preccurlyeq T$ and $|T|=|A|+1$.
The vertex and edge equations in
\noderef{TreeSequenceInitialSegment} therefore identify the unique last
addition exactly:
\[
                     T=A\cup\{e_1,y_1\}.                     \tag{1}
\]
Here $e_1$ is the actual edge entry of $T$ at index $|A|$, and $y_1$
is the actual new vertex entry.  Since $\ell-1=0$, the displayed
\noderef{TreeSequencePrefix} has length $|A|$; entrywise prefix equality
identifies it with $A$.  The assumed penultimate-prefix conclusion is
therefore
\[
              V(A)\text{ is elementary under }\varphi_n.     \tag{2}
\]

The `notElementary' field of
\noderef{HierarchyCounterexample}, after expanding
\noderef{ColoredVertexSetElementary}, supplies distinct vertices
$u,v\in V(T)$ and a color $\alpha$ missing at both.  By
\noderef{TreeSequenceVertexSet} and (1),
$V(T)=V(A)\cup\{y_1\}$.  Elementarity (2) prevents both vertices from
lying in $A$, and their distinctness prevents both from being $y_1$.
Interchanging their names if necessary and expanding
\noderef{TreeMissingColors} and \noderef{MissingColors} gives
\[
 \alpha\in\overline\varphi_n(A)
       \cap\overline\varphi_n(y_1).                           \tag{3}
\]

Let $n$ be the last positive series index.  Positivity in
\noderef{HierarchyCounterexample} rules out the initial constructor at
this index.  Eliminating the certificate from
\noderef{TashkinovSeries} and \noderef{TashkinovSeriesStep} leaves
exactly the SE, RE, and PE cases.

First suppose the final kind is SE.  The nonparallel field of
\noderef{HierarchyBase} and $q=0$ identify $A=T_n$.  An SE step first
adjoins its exact connecting edge $f_n$ to $T_n$ in the sense of
\noderef{TreeSequenceAdjoinEdge} and then takes a
\noderef{TashkinovClosure}.  The exact first-edge conclusion of
\noderef{HierarchyNonparallelFirstEdge}, specialized to the first
hierarchy stage $T$, identifies the edge at index $|T_n|$ with $f_n$.
Comparison with (1), including the adjoin certificate's outside
endpoint, yields
\[
                         f_n=e_1,\qquad u_n=y_1.               \tag{4}
\]
Thus this exact edge is the stored SE seed, not an ordinary later
augmentation.  The \noderef{MaximumDefectiveChoice} in the SE
constructor supplies its selected coloring.  The constructor's
`currentColoring' equation and the generating-coloring equation of
\noderef{ExtendedTashkinovTree} identify that coloring with
$\varphi_n$.  Reflexivity supplies the two clauses of
\noderef{StableColoring}, so the SE `seriesAlternative' says that the
colors missing at $u_n$ and those missing on $T_n$ are disjoint.
Equations (3)--(4) put $\alpha$ in both sets, a contradiction.  Hence
the final kind is not SE.

Suppose next that the final kind is RE, and let
$(\gamma_n,\delta_n)$ and the exact two-color cycle $O$ be its
\noderef{RevisitingExtensionCertificate}.  Again the nonparallel base
field gives $A=T_n$, and the same exact-entry argument gives (4).
The revisiting certificate chooses an earlier positive PE index $h$,
identifies its two stored colors with $\gamma_n,\delta_n$, and says that
every intervening index is RE.  At $h$, the
\noderef{ParallelExtensionColor} makes $\gamma_h$ missing at the PE
support under the selected coloring.  The exact defective edge in
\noderef{MaximumDefectiveChoice} is incident there and has color
$\delta_h$; the enlarged \noderef{StableColoring} certificate preserves
that edge and the missingness of $\gamma_h$.  Consequently the PE
\noderef{TwoColorComponentSwap} through the support replaces the
incident $\delta_h$ edge by a $\gamma_h$ edge and makes $\delta_h$
missing at that support in the PE output coloring.  Each intervening RE
constructor identifies its current coloring with its preceding one, and
the nested ladder trees retain the same support vertex.  Induction over
those exact constructor equations, followed by
$\delta_h=\delta_n$, proves
\[
                         \delta_n\in\overline\varphi_n(T_n).  \tag{5}
\]
The terminal RE `previousClosed' field is a
\noderef{ColoredVertexSetClosed} assertion on $T_n$, and its
`currentColoring' equation identifies that coloring with $\varphi_n$.
Thus (3), $A=T_n$, and (5) make both $\alpha$ and $\delta_n$ closed on
$T_n$.

If $\alpha=\delta_n$, cycle $O$ is already impossible: its exact edge
$f_n=e_1$ is incident with $y_1$, and a
\noderef{TwoColorCycle} has degree two at every incidence-graph node,
whereas $\delta_n=\alpha$ is missing at $y_1$ by (3).  Hence assume
$\alpha\ne\delta_n$.  Let $Q$ be the full
$(\delta_n,\alpha)$-component through $y_1$.  If $Q$ met $T_n$, a
first entry would be through an exact boundary edge colored
$\delta_n$ or $\alpha$, contradicting the two closedness assertions.
Therefore $Q$ is disjoint from $T_n$.  Since $y_1$ misses $\alpha$,
properness and \noderef{TwoColorIncidenceDegreeAtMostTwo} show that $Q$
is a possibly singleton \noderef{TwoColorPathComponent}; the singleton
case is retained.

Set
\[
 \sigma_n=\noderef{TwoColorComponentSwap}
       (G,\varphi_n,\delta_n,\alpha,\operatorname{inl}y_1).
\]
The component is disjoint from the component through the first
$\delta_n$-missing vertex of $T_n$.  The full-tail penultimate prefix
is exactly $A=T_n$, is elementary by (2), and both colors are missing
there by (3) and (5).  These are precisely the zero-divider Case 2
hypotheses of \noderef{StableChange}.  That lemma gives properness,
the same unique uncolored edge, and weak stability on $A$; expanding
\noderef{HierarchyWeaklyStable} in the non-PE case gives the exact
$(T_n,D_n,\varphi_n)$-\noderef{StableColoring} relation.  Applying
\noderef{Td} with the stated lower-rung package yields
\[
                         \noderef{ModColoring}(T,\sigma_n).    \tag{6}
\]
The exact one-edge seed in (4) is already the whole tree $T$, so the
empty augmentation sequence after that seed satisfies the RE
eligibility clause of (6).  Its universal conclusion provides an
actual \noderef{CorrespondingExtendedTashkinovTree} on the same ordered
$T$.  Extension-kind and exact connecting-edge-set agreement make its
terminal step RE with connecting set $\{e_1\}$.  The replacement
\noderef{RevisitingExtensionCertificate} therefore supplies a
two-color cycle through this same edge; equality of the raw stored
pairs identifies its unordered color pair with
$\{\gamma_n,\delta_n\}$.  The outside swap makes $\delta_n$ missing at
$y_1$, while the replacement cycle contains $e_1$ and hence $y_1$.
This contradicts the degree-two condition in
\noderef{TwoColorCycle}.  Hence the final kind is not RE.

It remains to suppose that the final kind is PE.  The parallel field of
\noderef{HierarchyBase} identifies
\[
                         A=T_n\vee R_n                        \tag{7}
\]
as the exact ordered union with the stored closure.  There is no SE/RE
seed exception in this case.  The final \noderef{TashkinovClosure} is
obtained from $T_n$ by a
\noderef{TashkinovAugmentationSequence}; its one-at-a-time length
increments pass through $|A|$ and $|T|=|A|+1$.  The intervening
\noderef{TashkinovAugmentation} is therefore the exact pair
$(e_1,y_1)$.  Put $\varphi_n(e_1)=\beta$.  Its augmentation witness
gives $\beta\in\overline\varphi_n(A)$, and incidence with $y_1$ plus
(3) gives $\alpha\ne\beta$.

The `parallelBase' field of \noderef{GoodHierarchy} gives at least four
colors in
\[
 (\overline\varphi_n(T_n)\cap\overline\varphi_n(R_n))
       -\bigcup_{\eta\in D_{n,0}}\Gamma^0_\eta.
\]
Choose $\theta$ there distinct from $\alpha$ and $\beta$.  We verify
the elementarity premise needed for its closedness without using a
zero-rung ETT.  If $n=1$, assertion (i) of
\noderef{TechnicalTheoremBase} applies to the exact rounds-one data.  If
$n>1$, \noderef{CertifiedETTRestriction} at rung $n-1$ gives an actual
positive lower-rung ETT displayed on $T_n$, with generating coloring
$\varphi_{n-1}$ and maximum property; lower assertion (i) applies to
that certified ETT.  Thus in both cases $T_n$ is elementary under the
pre-step coloring, and \noderef{ParallelCommonMissingClosed} makes
$\theta$ closed on the union (7):
\[
 \theta\in\overline\varphi_n(A),\qquad
 \partial_{\varphi_n,\theta}(A)=\varnothing.                 \tag{8}
\]

Apply conclusion (iii) of \noderef{interchange} to $\theta,\alpha$
under $\varphi_n$.  Properness and unique-uncoloredness are hypotheses;
the two factors of \noderef{CombinedStableColoring} are reflexive, and
(8) supplies missingness and closedness.  Hence the colors are
$A$-interchangeable.  Choose vertices of $A$ missing $\theta$ and
$\alpha$.  If they coincide, their full two-color component is the
isolated vertex component.  If they differ, properness,
\noderef{TwoColorIncidenceDegreeAtMostTwo}, and
\noderef{TwoColorComponentSpanningPath} make them the two ends of that
component.  Consequently the possibly singleton
$(\theta,\alpha)$-component $Q$ through $y_1$ cannot meet $A$: equality
with the component just described would either equate a nontrivial
component with an isolated vertex or add the third end $y_1$.

Swap $Q$ and call the result $\sigma_n$.  The exact penultimate prefix
is $A$ and is elementary by (2); both colors are missing there,
$\theta$ is missing on $T_n$, and $Q$ is disjoint from the component
through the first $\theta$-missing vertex.  Thus the zero-divider Case
2 clause of \noderef{StableChange} applies.  Its weak-stability
conclusion expands via \noderef{HierarchyWeaklyStable} to the two
factors of \noderef{CombinedStableColoring} on $T_n$ and $R_n$.

We also record the needed pointwise facts directly.  Since $Q$ has no
vertex in $A$, no edge-node incident with $A$ can lie in $Q$, by the
connected-support definition of \noderef{TwoColorComponent}.  Hence
the cases in \noderef{TwoColorComponentSwap} give
\[
 \sigma_n(f)=\varphi_n(f)
 \quad\text{for every edge $f$ incident with $A$}.            \tag{9}
\]
Thus all vertex missing sets and all colored boundary edges of $A$ are
unchanged: $\theta$ remains missing and closed on $A$, and $\beta$
remains missing there.  Since $\beta$ differs from both swap colors,
the exact edge is fixed separately,
$\sigma_n(e_1)=\varphi_n(e_1)=\beta$, while $y_1$ misses $\theta$
after the swap, including when $Q$ was singleton.

Apply conclusion (iii) of \noderef{interchange} under $\sigma_n$ to
$\theta,\beta$.  The common conclusion of
\noderef{StableChange} supplies properness and unique-uncoloredness, its
weak-stability conclusion supplies the combined stability, and (9)
supplies missingness and closedness.  It follows that these colors are
$A$-interchangeable.  Their missing witnesses in $A$ determine either
an isolated component when they coincide or a spanning path with those
two ends when they differ.  The component containing the exact edge
$e_1$ also meets $A$ and has the additional end $y_1$, where $\theta$
is missing.  The two components are therefore distinct in both cases.
This contradicts \noderef{TreeInterchangeable}, whose path components
include isolated vertices.  The PE case is impossible as well.

All three possible terminal constructors have produced contradictions.
Therefore no zero-divider hierarchy counterexample satisfying the exact
one-addition and penultimate-elementarity hypotheses exists.
\end{proof}
